% Generated draft: requires mathematical and source review.
\documentclass[11pt]{article}
\usepackage[margin=25mm]{geometry}
\usepackage{fontspec}
\setmainfont{Latin Modern Roman}
\usepackage{amsmath,amssymb,mathtools,unicode-math}
\setmathfont{Latin Modern Math}
\usepackage{enumitem,needspace,xurl,hyperref}
\setlength{\parindent}{0pt}
\setlength{\parskip}{5pt}
\setlength{\emergencystretch}{3em}
\newcommand{\N}{\mathbb{N}}
\newcommand{\Z}{\mathbb{Z}}
\newcommand{\Q}{\mathbb{Q}}
\newcommand{\R}{\mathbb{R}}
\newcommand{\C}{\mathbb{C}}
\newcommand{\F}{\mathbb{F}}
\newcommand{\A}{\mathbb{A}}
\newcommand{\PP}{\mathbb{P}}
\newcommand{\E}{\mathbb{E}}
\newcommand{\eps}{\varepsilon}
\newcommand{\dd}{\,\mathrm{d}}
\begin{document}
\section*{UnsolvedMath ID: 20001274}
\textbf{Source problem code:} catalogue entry
\subsection{Definitions and mathematical statement}
Fix integers \(n\geq3\) and \(2\leq k\leq n-1\). A convex body is a compact convex subset of \(\R^n\) with nonempty interior. Write \(G(n,k)\) for the Grassmannian of linear \(k\)-subspaces and \(K\mid S\) for orthogonal projection. The restored source question asks whether
\[
 \forall S\in G(n,k)\ \exists Q_S\in O(S),\ t_S\in S:
 K\mid S=Q_S(L\mid S)+t_S
\]
implies
\[
 \exists\sigma\in\{1,-1\},\ t\in\R^n:\quad K=\sigma L+t.
\]
The local rigid motions may depend on \(S\) and may reverse orientation. The global alternatives are translation of \(L\) or its reflection in the origin \(-L\), not arbitrary global orthogonal congruence. The variable \(S\) missing before \(G(n,k)\) in the extraction is restored using the supplied verification.

The aid reports Ning Zhang's 2018 refutation in the hyperplane case: for every \(n\geq3\), there are convex bodies \(K,L\subseteq\R^n\) for which no global \(\sigma,t\) as above exist, but for every \(S\in G(n,n-1)\) an origin-fixing \(Q_S\in O(S)\) satisfies \(K\mid S=Q_S(L\mid S)\). This is a counterexample to the general proposed implication; the report does not separately settle every lower-dimensional \(k\) in a fixed ambient dimension.

The aid also reports a positive special case. Define the difference body \(D_K=K-K\), with no factor of \(1/2\), and for an origin-symmetric projected body define
\[
 \operatorname{Sym}_S(D\mid S)=\{Q\in O(S):Q(D\mid S)=D\mid S\}.
\]
Under the local congruence hypothesis, suppose additionally
\[
 D_K=D_L=:D,\qquad
 \operatorname{Sym}_S(D\mid S)=\{I_S,-I_S\}\quad\text{for every }S\in G(n,k).
\]
Then the reported conclusion is \(K=L+t\) or \(K=-L+t\) for some \(t\in\R^n\). In support-function notation \(h_K(u)=\sup_{x\in K}\langle x,u\rangle\), the aid describes the reduction to
\[
 h_K(u)=h_L(\sigma_Su)+\langle t_S,u\rangle\quad(u\in S),\qquad \sigma_S\in\{1,-1\},
\]
and invokes Myroshnychenko's established sign-and-linear-term gluing theorem to obtain one global sign and linear term. Equivalently, \emph{assuming} equality of difference bodies and local congruences, every counterexample must have at least one \(S\) whose projected difference body has an orthogonal symmetry beyond \(\pm I_S\). Equality of difference bodies is an extra hypothesis and is not asserted to follow from arbitrary projectionwise congruences.

The original notes report an affirmative counterpart for \(k=2\) and rotations in reference [R], by Ryabogin, and recent progress by F. Nazarov communicated by Ryabogin. Those source-time qualifications do not supply the exact hypotheses of the rotation result or a separate theorem of Nazarov. The appended ``Vertex Index Problems, Proposed by Alexander Litvak'' is the heading of the next section, not attribution for this question.
\subsection{Short English statement}
The question compares congruent projections in every fixed-dimensional subspace with global equality up to translation and central reflection. The aid reports hyperplane counterexamples and a positive special case requiring equal difference bodies with no projected orthogonal symmetries beyond the two signs.
\subsection{Sources}
\begin{itemize}
\item Local registry: permanent UnsolvedMath ID \(20001274\), exact public code \texttt{catalogue entry}; \texttt{aim-convex-geometry-notes.json}, \texttt{notes[5]}.
\item AIM workshop ``Sections of convex bodies'', Section 8, problem 5: \url{https://aimath.org/pastworkshops/sectionsconvexproblems.pdf}. The source points to [G, Problem 3.2 and Note 3.1] and Ryabogin [R]. Full entries for [G] and [R] are not supplied. The source reports F. Nazarov's progress as communicated by Ryabogin, printed once as ``Dimtry Ryabogin''.
\item The supplied verification restores \(S\in G(n,k)\), the projection notation and the range \(2\leq k\leq n-1\), and identifies the following Litvak heading as spillover. The registry marks the statement \texttt{reconstructed\_unverified}.
\item Research attempt 1: \texttt{counterexample}, status at run \texttt{refuted}, rigor label \texttt{proved}, novelty confidence low, classification \texttt{SOLVED-IN-LITERATURE}. It attributes the hyperplane counterexamples to Ning Zhang (2018) and the gluing theorem to Myroshnychenko, without full bibliographic entries, and describes the difference-body symmetry criterion as a candidate special case. The current registry status is solved.
\end{itemize}
\end{document}
